% =====================================================================
%  commands.tex
%
%  Custom commands.  The most important one is \sectag: a small margin
%  label that flags a section as Core / Bridge / Proof / Counterexample
%  / Computing / Optional / Project, matching the TOC's tag system.
% =====================================================================

% --- Section tags (Core, Bridge, Proof, ...) -------------------------
% Usage: \section{Title}\sectag{Core}
% Or to combine: \section{Title}\sectag{Core, Proof}
%
% We use \marginnote (not \marginpar) so the tag stays anchored next
% to the section title even on a page with little body text.  The
% optional [vertical] offset pulls the tag up by ~3 baselines so it
% lines up with the section heading rather than the first subsection.
\newcommand{\sectag}[1]{%
  \marginnote{%
    \footnotesize\sffamily\color{Sepia}%
    \fcolorbox{Sepia}{white}{%
      \begin{minipage}{\dimexpr\marginparwidth-2\fboxsep-2\fboxrule\relax}%
        \centering\strut #1\strut%
      \end{minipage}%
    }%
  }[-3\baselineskip]%
}

% Shorter aliases for the seven tag types — purely a convenience.
\newcommand{\Core}{\sectag{Core}}
\newcommand{\Bridge}{\sectag{Bridge}}
\newcommand{\Proof}{\sectag{Proof}}
\newcommand{\Counter}{\sectag{Counterexample}}
\newcommand{\Computing}{\sectag{Computing}}
\newcommand{\Optional}{\sectag{Optional}}
\newcommand{\Project}{\sectag{Project}}

% Combine two tags neatly, e.g. \CoreProof, \CoreBridge, etc.
\newcommand{\CoreProof}{\sectag{Core / Proof}}
\newcommand{\CoreBridge}{\sectag{Core / Bridge}}
\newcommand{\CoreComputing}{\sectag{Core / Computing}}
\newcommand{\OptionalProject}{\sectag{Optional / Project}}
\newcommand{\OptionalBridge}{\sectag{Optional / Bridge}}
\newcommand{\ComputingProject}{\sectag{Computing / Project}}
\newcommand{\ProofOptional}{\sectag{Proof / Optional}}

% --- Pedagogical inline markers --------------------------------------
% A small "hook" callout the chapter ends use to tease the next chapter.
\newcommand{\hook}[1]{%
  \par\medskip
  \noindent\textbf{\color{NavyBlue}Hook.}\quad #1\par\medskip
}

% Margin "tip" / "note" used sparingly inside body text.
\newcommand{\margintip}[1]{\marginpar{\footnotesize\itshape #1}}

% --- Common math operators (declared once, used everywhere) ----------
\DeclareMathOperator{\sgn}{sgn}
\DeclareMathOperator{\dom}{dom}
\DeclareMathOperator{\range}{range}
\DeclareMathOperator{\arsinh}{arsinh}
\DeclareMathOperator{\arcosh}{arcosh}
\DeclareMathOperator{\artanh}{artanh}

% --- Number sets ------------------------------------------------------
\newcommand{\R}{\mathbb{R}}
\newcommand{\Q}{\mathbb{Q}}
\newcommand{\Z}{\mathbb{Z}}
\newcommand{\N}{\mathbb{N}}
\newcommand{\C}{\mathbb{C}}

% --- Calculus shortcuts ----------------------------------------------
% Differential d (upright): use as \diff x, \diff t, etc.
\newcommand{\diff}{\mathop{}\!\mathrm{d}}

% Derivative as a fraction: \dd{y}{x}, \ddn{2}{y}{x}
\newcommand{\dd}[2]{\dfrac{\diff #1}{\diff #2}}
\newcommand{\ddn}[3]{\dfrac{\diff^{#1} #2}{\diff #3^{#1}}}

% Partial derivative: \pd{f}{x}, \pdn{2}{f}{x}
\newcommand{\pd}[2]{\dfrac{\partial #1}{\partial #2}}
\newcommand{\pdn}[3]{\dfrac{\partial^{#1} #2}{\partial #3^{#1}}}

% Definite integral with neat spacing: \dint{a}{b}{f(x)}{x}
\newcommand{\dint}[4]{\int_{#1}^{#2}\! #3 \,\diff #4}

% Limit shortcut: \lmt{x}{a}{ f(x) }
\newcommand{\lmt}[3]{\lim_{#1\to #2} #3}

% --- Vectors & norms -------------------------------------------------
\newcommand{\vect}[1]{\mathbf{#1}}
\newcommand{\abs}[1]{\left\lvert #1 \right\rvert}
\newcommand{\norm}[1]{\left\lVert #1 \right\rVert}

% --- Convenience environments ----------------------------------------
% A "Project" environment for the chapter-end project subsections.
\newenvironment{project}[1][Project]%
  {\par\medskip\noindent\textbf{\color{ForestGreen}#1.}\quad}%
  {\par\medskip}

% A "Discovery" callout for guided-discovery problems.
\newenvironment{discovery}[1][Discovery]%
  {\par\medskip\noindent\textbf{\color{Plum}#1.}\quad}%
  {\par\medskip}

% A simple "Worked example" environment whose body is set off in a box.
% (Optional; remove if you prefer plain example numbering.)
\newenvironment{worked}[1][Worked Example]%
  {\par\medskip\noindent\textsc{#1.}\quad\itshape}%
  {\par\medskip}
